\documentclass{article}
\usepackage{graphicx} % Required for inserting images
\usepackage[document]{ragged2e}
\title{R\&D\_cs691\_progress\_report}
\author{Rajat Ranjan Meher(25m0777) \\[0.2cm]
\textbf{Instructor:-} Bhaskaran Raman\\
[0.2cm]}
% \vspace{20cm}
\date{Even semster : 2026}


\begin{document}

\maketitle

\section{1st week (13th Jan - 20th Jan) (12 Hours)}

\subsection{Objective:-Learn about Androids in Java} 
\vspace{0.2cm}

\textbf{Android terms:-}\\

Activity,View,layout,Intent,Event,Event listener(onClick),Debugging using (Log.d),Context ,Type  of text , use  of xml, .java and AndroidManifest\par
\vspace{2mm}
\textbf{some methods which i remember:-}\par
\vspace{2mm}
R.id.idname ,findViewById(),setTitle(), getStringExtra(),putExtra() \par
\vspace{2mm}
\textbf{Application of concepts:-}\par
\vspace{2mm}
Going from one activity to another activity onClick of button using intent.\\
Using \textbf{toast} for pop up \\
Sending text between different activities , using Intent and putExtra , get string by Extra.\\
Using parentactivity in androidManifest for arrow to main activity.\\
Multiple text line from password, autofill, multiline, single line\\
Can also be used interaction between different view inside a same activity.
\par
\vspace{2mm}
\textbf{Some learnings:-}\par
\vspace{2mm}
You have to typecast findViewById to either button , EditView etc for any application of methods on it. \\
android:lable in mainactivity\_xml change the label of main activity and also the app's name. So for just app's name use label in main activity , while for main activity use setTitle() in the oncreate in mainactivity.java\\
For same string to use in different place use name and value inside values>themes.xml
\section{2nd week (20th Jan - 27th Jan) (10 Hours)}
\subsection{Objective:-Learn about Androids in Java} 
\textbf{Language:-}Familiarization with Kotlin syntax. As many of the file are written .kt.\par
\vspace{2mm}
\textbf{Views explore:-} Explored many view such as layout and text with many features.\par
\vspace{2mm}
\textbf{Library documentation:-} Read documentation like Intent.java, Activity.java , Toast.java. \par
\vspace{2mm}
\textbf{Mistakes and Learnings:-} Using image above a certain size causes crash in app. Solved it by taking a screenshot to use it.
\section{3nd week (5th Feb - 12th Feb) (10 Hours)}
\subsection{Objective:-Learn About Androids in Java and Kotlin}
\textbf{Summary writing:-}Made summary of CS 490 RND project. \par
\vspace{2mm}
\textbf{Language:- } More practise on Kotlin to get more familiarity. \par
\vspace{2mm}
\textbf{Architecture:-}Read a bit about MVI and started exploring MVVM.
\par \vspace{2mm}
\textbf{New Terms :-} Singleton , LiveData , MODEL, VIEW and VIEWMODEl
\par \vspace{2mm}

\section{4th week (1st March - 7th March) (10 Hours)}
\subsection{Objective:-Understanding Project Architecture for VoIP Measurement}
\textbf{Project Understanding:-} Studied modified problem statement focusing on Voice-over-IP (VoIP) calls such as WhatsApp and Google Meet.\par
\vspace{2mm}
\textbf{Architecture:-} Learned Client-Server architecture where Android app periodically collects network data and communicates with Python backend.\par
\vspace{2mm}
\textbf{Modules Identified:-}
\begin{itemize}
\item Periodic Measurement Module (Timer-based triggering)
\item Network Sensing Module (WiFi and Mobile data collection)
\item Backend Aggregation Module
\end{itemize}
\vspace{2mm}
\textbf{Android Concepts:-} Studied ConnectivityManager, WifiManager and background execution.\par
\vspace{2mm}
\textbf{Learnings:-} Understood that VoIP calls cannot be detected using Telephony APIs and require indirect measurement through network conditions.

\section{5th week (8th March - 14th March) (10 Hours)}
\subsection{Objective:-Implementation of Periodic Network Monitoring}
\textbf{Coding Work:-}
\begin{itemize}
\item Implemented periodic task using Handler to trigger data collection
\item Collected active network type (WiFi or Mobile) using ConnectivityManager
\item Handled cases where no network is available
\end{itemize}
\vspace{2mm}
\textbf{Manifest Work:-}
\begin{itemize}
\item Added permissions (ACCESS\_NETWORK\_STATE, ACCESS\_WIFI\_STATE)
\item Configured app for background data collection
\end{itemize}
\vspace{2mm}
\textbf{Learnings:-} Understood shift from event-driven system to time-based measurement for VoIP scenarios.

\section{6th week (15th March - 21st March) (10 Hours)}
\subsection{Objective:-Network Data Collection and User Feedback}
\textbf{UI Design:-}
\begin{itemize}
\item Designed simple feedback interface for user-reported call quality
\item Added RadioGroup for quality (Good, Fair, Poor)
\item Included basic issue selection using CheckBoxes
\end{itemize}
\vspace{2mm}
\textbf{Data Collection:-}
\begin{itemize}
\item Extracted WiFi signal strength (RSSI) using WifiManager
\item Identified network type (WiFi / Mobile)
\item Combined user feedback with network parameters
\end{itemize}
\vspace{2mm}
\textbf{Learnings:-} Learned how to correlate user experience with underlying network conditions in VoIP calls.

\section{7th week (22nd March - 31st March) (10 Hours)}
\subsection{Objective:-Backend Integration and Data Transmission}
\textbf{Backend Work:-}
\begin{itemize}
\item Studied Flask API for receiving periodic measurements
\item Created structured JSON including quality, network type and signal strength
\item Implemented HTTP POST request using OkHttp
\end{itemize}
\vspace{2mm}
\textbf{System Integration:-}
\begin{itemize}
\item Integrated periodic measurement with data transmission
\item Ensured data is sent asynchronously without blocking UI
\end{itemize}
\vspace{2mm}
\textbf{Learnings:-} Understood full pipeline of VoIP performance measurement using indirect network monitoring.
\end{document}
